\chapter{算法与编程}
\section{数据摘要与区间扫描}
\noindent\textbf{要点与方法.}\enspace
处理数据流时，先确定需要保留的摘要。均值与方差可由样本量、均值和中心平方和共同维护；合并时还要计入组间均值差。对连续区间的极值，单调队列只保留尚未过期且未被更晚元素支配的候选。若扩大窗口只会使约束更难满足，便可通过单向移动左端点寻找最长可行区间。

正确性依靠摘要或队列的不变量，线性时间则依靠每个元素至多入队、出队一次。随机模拟另需证明接受后的分布与终止性；拒绝后重新抽样时，必须使用新的独立试验。
\par\smallskip


\problem{A31}{分组摘要的合并与拆分}{}
两组互不重叠的数据只保存摘要 $(n_A,\mu_A,M_A)$ 与 $(n_B,\mu_B,M_B)$，其中 $n_A,n_B\ge1$，$M_A=\sum_{i\in A}(x_i-\mu_A)^2$，$M_B$ 同理，原始观测不可读取。求合并后的均值、中心平方和及两种方差。若已知合并摘要和 A 组摘要，怎样恢复 B 组摘要？讨论运算量及数值风险。
\solution 记 $n=n_A+n_B$、$\delta=\mu_B-\mu_A$。总和相加立即给出
\[
\mu=\mu_A+\frac{n_B}{n}\delta.
\]
对 A 组，将 $x_i-\mu=(x_i-\mu_A)+(\mu_A-\mu)$ 展开平方并求和。组内偏差和为零，故交叉项消失；B 组同理。因此
\[
M=M_A+M_B+n_A(\mu_A-\mu)^2+n_B(\mu_B-\mu)^2
 =M_A+M_B+\frac{n_An_B}{n}\delta^2.
\]
总体方差为 $M/n$，分母作自由度修正的样本方差为 $M/(n-1)$。后者的无偏性还需要相应的独立同分布抽样假设。

反向拆分时，先取 $n_B=n-n_A>0$，然后
\[
\mu_B=\frac{n\mu-n_A\mu_A}{n_B},\qquad
M_B=M-M_A-\frac{n_An_B}{n}(\mu_B-\mu_A)^2.
\]
这些恒等式只需常数次运算和常数大小的摘要。精确算术下可以按任意二叉合并顺序得到同一结果，因为每一步都恢复真实的总均值和中心平方和；浮点舍入会使不同顺序略有差异。逆向拆分可能发生严重相消，特别是 B 组很小而两项被减数很大时。明显的负 $M_B$ 应视为摘要不一致或数值失败，不能一概截成零。
\insight 总方差由组内波动和组间均值差共同组成；直接平均两组方差会漏掉最后一项。

\problem{A32}{波动受限的最长连续区间}{}
给定实数序列 $x_0,\ldots,x_{T-1}$ 和阈值 $D\ge0$，找出满足
\[
\max_{i\le t\le j}x_t-\min_{i\le t\le j}x_t\le D
\]
的最长非空连续区间。长度相同时返回起点最早者；空序列返回空区间。设计 $O(T)$ 时间算法，并证明左端点的移动不会跳过最优解。
\solution 维护左端点 $\ell$、一个存放递减数值下标的最大值队列，以及一个存放递增数值下标的最小值队列。依次加入右端点 $r$：在最大队列尾部删除值不大于 $x_r$ 的下标，在最小队列尾部删除值不小于 $x_r$ 的下标，再将 $r$ 入队。

只要两队首对应的差超过 $D$，就将 $\ell$ 右移一位，并删除小于新 $\ell$ 的队首下标。停止时两队首就是窗口的真实极值。记录长度 $r-\ell+1$，仅在严格超过当前最好长度时更新答案；由于右端点递增，这保留了最早起点。

正确性的关键是：若 $[\ell,r]$ 不合法，则任何包含它的区间也不合法，故把这个左端点丢弃不会失去以当前或未来位置结尾的可行解。反之，算法停止收缩时，$\ell$ 是当前右端点下仍可能合法的最早起点，因此得到以 $r$ 结尾的最长合法区间。遍历所有 $r$ 即得到全局最优。

队尾被删除的下标被一个更晚且更极端的数支配，既不可能先于新数继续留在窗口，也不会成为需要保留的极值。每个下标至多入队、出队各一次，左端点也至多移动 $T$ 次，因此总时间 $O(T)$，最坏空间 $O(T)$。非空输入中，每个单点区间都可行。

\problem{P09}{用未知偏置硬币模拟骰子}{}
只能独立抛掷同一枚硬币，其未知正面概率为 $p\in(0,1)$，没有额外随机源。设计一个不需要知道 $p$、能够精确输出 $1,\ldots,6$ 中等概率一个数的算法。证明输出均匀且几乎必然结束，并求你所给算法的期望抛掷次数。
\solution 先产生公平比特。连续抛两次；出现 HT 时输出0，出现 TH 时输出1，HH或TT则丢弃并重试。HT和TH的概率同为 $p(1-p)$，故条件于接受时输出公平。不同输出使用不相交的独立试验段，下一段在前一段结束后仍有相同的分布，因此输出比特独立。

设生成一个公平比特所需抛掷次数的均值为 $b$。每对接受概率为 $2p(1-p)>0$，重试次数有几何尾，且
\[
b=2+\bigl(1-2p(1-p)\bigr)b
  =\frac1{p(1-p)}.
\]
每次取三个公平比特，读为整数 $Z\in\{0,\ldots,7\}$。若 $Z\le5$ 则输出 $Z+1$，否则丢弃整组三比特并重试。每个整数等概率，接受后的六个结果仍等概率，且每轮接受概率为 $3/4$。

令总抛掷次数均值为 $E$。一轮平均使用 $3b$ 次硬币抛掷；被拒绝后使用全新的独立试验。因此
\[
E=3b+\frac14E,\qquad
\boxed{E=\frac4{p(1-p)}}.
\]
两层接受概率均为正，故算法几乎必然结束。这里求的是这套具体算法的成本，并未声称它是最优算法；当 $p$ 接近0或1时，提取公平比特的代价会很大。


\par\Needspace{8\baselineskip}
\subsection*{小节练习}
\exercise{OA31}{滚动均值与方差}
实数数据流 $x_1,x_2,\ldots$ 按顺序到达。给定窗口长度 $N\ge1$，每次新数据到来时，更新最近 $N$ 个数据的均值与方差，要求更新时间尽量低。数据不足 $N$ 个时使用所有已到达数据。这里方差按窗口内平方偏差的平均定义。给出算法、更新公式及时间与空间复杂度。

\exercise{OA32}{滚动最大值}
给定按时间顺序到达的实数数据流 $x_1,x_2,\ldots$ 和正整数 $N$，每次新数据到来时输出最近 $N$ 个数据的最大值；不足 $N$ 个时使用所有已有数据。设计高效算法，证明其正确性，并区分单次最坏时间复杂度与摊还时间复杂度。


\section{回撤区间}
\noindent\textbf{要点与方法.}\enspace
对正价格序列，先固定终点 $j$。回撤率 $1-p_j/p_i$ 随起点价格 $p_i$ 增大而增大，因此最大回撤可化为查询合法起点中的最高价；回撤持续时间的限制决定候选窗口，价格必须严格为正。若目标改为达到给定回撤的最短时间，则应优先保留更晚且价格更高的起点，不能直接沿用最大回撤的答案。

删除候选前要说明它为何不可能改善未来结果：或已过期，或被另一个候选支配，或已取得自己的最短可行区间。查询与入队的先后顺序保证 $i<j$；相等价格的保留规则和更新条件共同决定并列时选哪个区间。
\par\smallskip


\problem{A33}{限定持续时间的最大回撤}{}
给定严格为正的价格 $p_0,\ldots,p_{T-1}$ 和整数 $L\ge1$。只允许持续 $1$ 至 $L$ 个时间单位的回撤区间，即 $1\le j-i\le L$。求
\[
D=\max\left(0,\max_{1\le j-i\le L}\left(1-\frac{p_j}{p_i}\right)\right).
\]
若 $T<2$，约定 $D=0$。若 $D>0$，返回一个最优区间，先按终点、再按起点从早到晚打破平局；否则返回零和空区间。给出线性时间算法。
\solution 对固定终点 $j$，因为 $p_j>0$，回撤率随分母 $p_i$ 增大而增大。因此只须找合法前缀窗口
\[
\max(0,j-L)\le i\le j-1
\]
内的最高价格。用递减单调队列维护这些候选：在处理 $j$ 前，加入 $j-1$，删除队尾价格严格小于 $p_{j-1}$ 的项；相等价格保留，以便选较早起点。随后删除下标小于 $j-L$ 的队首。

若队首为 $i$，计算 $1-p_j/p_i$。从 $j=1$ 到 $T-1$ 扫描，只有严格大于当前最好正回撤时才更新。这样较早终点优先，而队内相等峰值保证较早起点优先。$T<2$ 时直接返回零和空区间。

被队尾删除的价格被更晚且更高的价格支配：后者在任何未来窗口内至少同样有利，且不会更早过期。队首因而始终是合法窗口中的最高价格；逐终点取最优再取全局最优证明算法正确。每个下标进出队列至多一次，时间 $O(T)$，空间 $O(\min(T,L))$。
\pitfall 全历史最高价可能已超过允许的持续时间，直接用累计最大值会得到不可行区间。

\problem{A34}{首次达到回撤阈值的最短区间}{}
给定严格为正的价格序列及 $d\in(0,1)$，寻找最短区间 $(i,j)$，满足 $i<j$ 且 $p_j\le(1-d)p_i$。若不存在则报告无解；长度并列时选最早终点。设计 $O(T)$ 时间算法。说明为什么寻找最大回撤本身不能解决这个问题。
\solution 维护价格严格递减的候选起点队列。处理终点 $j$ 时，先反复检查队首 $i$：若
\[
p_j\le(1-d)p_i,
\]
就用 $j-i$ 更新最短长度，然后弹出 $i$；否则停止这轮检查。随后删除队尾价格不大于 $p_j$ 的下标，再把 $j$ 入队。先查询后入队保证 $i<j$。

队首不达阈值时，后面价格更低的候选也不达阈值，故可以停止。达阈值的起点被弹出是安全的：对固定 $i$，当前 $j$ 是算法尚未处理过的最早可用终点，任何未来终点都会产生更长的区间。若有更晚起点价格不低于一个旧起点，则新起点在未来达到相同阈值只会更容易，区间也更短，因此旧起点可以从队尾删除。

这两种删除分别由“已经得到该起点的最好长度”和“被更晚、更高的价格支配”保证。每个候选至多进出队列一次，时间 $O(T)$，空间 $O(T)$。按终点顺序扫描并仅在长度严格缩短时更新，可保留最早终点。最大回撤优化的是下降比例，可能需要很长时间，与这里的最短持续时间是不同目标。

\par\Needspace{12\baselineskip}
\subsection*{小节练习}
\exercise{OA33}{最大回撤}
给定正价格序列 $p_1,\ldots,p_n$，$n\ge1$。最大回撤定义为从较早价格到较晚价格的最大相对跌幅：
\[
 D=\max\left(0,\max_{1\le i<j\le n}\frac{p_i-p_j}{p_i}\right),
\]
其中 $n=1$ 时取 $D=0$。给出 $O(n)$ 时间算法并证明正确性，说明空间复杂度。若 $D>0$，如何同时返回一个实现最大回撤的起止下标？


\section{算法习题}
\noindent\textbf{要点与方法.}\enspace
动态规划先定义状态所包含的信息，再写转移与边界。删除次数、库存、已选区间数等限制必须进入状态；不可达状态应与合法的零收益状态区分。计数题还要保证各条转移对应互不重叠的方案，否则会重复计数。

选择算法时，应证明可利用的结构：二分依赖单调的排除方向，贪心依赖交换或支配关系，分治需要说明跨子问题的贡献如何计算。复杂度按所有状态、转移和数据结构操作累计；空间压缩时还须检查更新顺序是否误用本轮结果。随机算法应另行区分分布正确性、期望成本与最坏成本。
\par\smallskip


\algproblem{ALG01}{允许删除一次的最大子段和}
给定非空实数数组 $a_0,\ldots,a_{T-1}$。选取一个连续区间，可以从该区间删除至多一个元素，但至少保留一个。求保留元素之和的最大值，要求 $O(T)$ 时间、$O(1)$ 额外空间。
\solution 令 $f_i$ 为以 $i$ 结尾且不删除元素的最好值，$g_i$ 为以 $i$ 结尾、至多删除一个元素且仍非空的最好值；在 $g_i$ 中，位置 $i$ 本身可以被删。初始化 $f_0=g_0=a_0$。对 $i\ge1$，
\[
f_i=\max(a_i,f_{i-1}+a_i),\qquad
g_i=\max(a_i,g_{i-1}+a_i,f_{i-1}).
\]
最后一式依次对应重新开始、保留 $a_i$ 并延长已有区间、删除当前 $a_i$。若某个解删除了区间第一个元素，把区间起点后移即可得到相同保留序列，所以不需单独增加这种情况。其余合法解按末元素是否删除恰被递推覆盖。

答案为 $\max_i g_i$。更新时保存旧的 $f$，避免第三项误读新状态。始终以实际元素初始化，不能用零替代，否则全负数组会错误地返回空段。两状态的归纳证明同时保证最优性；时间 $O(T)$、额外空间 $O(1)$。

\algproblem{ALG02}{到达终点的最少跳跃次数}
非空数组 $a_0,\ldots,a_{T-1}$ 为非负整数。从位置0出发，在位置 $i$ 可向右跳 $1,\ldots,a_i$ 格。求到达 $T-1$ 的最少跳数；不能到达则返回无解。要求线性时间、常数额外空间，不能假设输入保证可达。
\solution $T=1$ 时返回0。维护当前层可达的最右位置 $R$、扫描到目前为止下一层可达的最右位置 $F$，以及已使用的层数 $k$，初始均为0。从 $i=0$ 扫到 $T-2$，先置
\[
F\leftarrow\max\{F,\min(T-1,i+a_i)\}.
\]
每当 $i=R$，当前层已全部扫描。若 $F=R<T-1$，则没有新的可达位置，返回无解；否则令 $k\leftarrow k+1$、$R\leftarrow F$。一旦 $R=T-1$ 就返回 $k$。

从一个位置允许跳到一整段前方位置，所以至多 $k$ 次跳跃能到达的位置构成从0开始的连续前缀。扫描完此前缀后，所有允许的一步延伸的并集恰到 $F$，因此新的边界是至多 $k+1$ 步的最远位置。算法逐层推进，首次覆盖终点的层数就是最短跳数。每个下标只扫描一次，时间 $O(T)$，额外空间 $O(1)$。

\algproblem{ALG03}{持有期限受限的一次交易}
价格 $p_0,\ldots,p_{T-1}$ 非负，给定整数 $1\le h\le L$。最多做一次交易，买入日 $i$ 与卖出日 $j$ 必须满足 $h\le j-i\le L$，每次卖出收取费用 $c\ge0$。可选择不交易。求最大净利润和一个最优日期对，要求 $O(T)$ 时间。
\solution 对卖出日 $j$，合法买入日构成
\[
\max(0,j-L)\le i\le j-h.
\]
若该区间为空，则当天没有可行交易。维护价格递增的单调队列：在第 $j$ 天把 $j-h$ 加入（若非负），入队时删去队尾价格严格更高的下标；再删除小于 $j-L$ 的队首。保留相等价格可选择较早买入日。

队首给出最低合法买价，候选利润为 $p_j-p_i-c$。答案从零开始，仅在严格改善时记录日期；若最终仍为零，返回不交易。若 $T\le h$，自然没有合法日期对。

固定卖出日时，费用与卖价固定，最低买价必定最优。队尾删除由更晚且更低的价格支配旧价格来保证；旧价格不会比新价格更晚过期。延迟加入保证最短持有期，队首过期保证最长持有期。每个下标只入队、出队一次，时间 $O(T)$，空间 $O(\min(T,L))$。

\algproblem{ALG04}{有限库存的逐日交易}
给定 $T$ 天的非负价格及整数库存上限 $K\ge0$。每天至多买入一单位、卖出一单位或不操作，禁止做空，买卖每单位均收手续费 $c\ge0$。初始现金和库存均为0，允许现金为负，最后库存必须为0。求最大终值，并给出 $O(T(K+1))$ 的动态规划。
\solution 令 $V_t(q)$ 为处理前 $t$ 天、持有 $q$ 单位时的最大现金，$0\le q\le K$。初值 $V_0(0)=0$，其余为 $-\infty$。在价格 $p_t$ 的一天，
\[
V_{t+1}(q)=\max\left\{
\begin{array}{l}
V_t(q),\\
V_t(q-1)-p_t-c\quad(q\ge1),\\
V_t(q+1)+p_t-c\quad(q<K).
\end{array}\right\}
\]
三项分别是不操作、买入、卖出。所有右端必须读取上一天的数组；若原地连续更新，可能错误地允许一天买卖多单位。

任意合法策略按最后一天的动作恰落入以上三类。去掉最后动作后得到相应旧状态，而每个有限旧状态接上该动作也合法，归纳证明递推同时给出上界与可达值。最终答案为 $V_T(0)\ge0$，因为始终不操作可行。时间 $O(T(K+1))$，滚动数组空间 $O(K+1)$；保存动作还可恢复交易日程。$K=0$ 或 $T=0$ 时答案都是0。允许现金为负是必要的模型说明，否则还需检查每笔买入是否有足够现金。

\algproblem{ALG05}{至多穿过一处障碍的路径计数}
一个 $R\times C$ 网格的每格标为0或1，1表示障碍。只能从左上角向右或向下移动到右下角，最多允许经过一个障碍格，起终点也按同一规则计入。求不同路径数对给定整数 $M\ge2$ 的余数，要求 $O(RC)$ 时间及 $O(C)$ 空间。
\solution 令 $A_{i,j}$、$B_{i,j}$ 分别计到达 $(i,j)$ 时经过零个、恰好一个障碍的路径数。先把上方和左方的两种计数分别相加，记为 $a,b$。若当前格为0，置 $(A,B)=(a,b)$；若为1，置 $(A,B)=(0,a)$。每次加法均取模 $M$。

起点单独初始化：起点为0则 $(A,B)=(1,0)$，为1则为 $(0,1)$。其余越界前驱一律贡献零。终点答案为 $A_{R-1,C-1}+B_{R-1,C-1}\pmod M$。

每条路径的倒数第二格恰为上方或左方，两类互斥；障碍计数又恰由当前格是否为1决定，所以递推既不遗漏也不重复。按行从左到右更新两个长度 $C$ 的数组，更新前的本列保存上方状态、已更新的前一列保存左方状态，先取出两者再覆盖即可。$R,C\ge1$ 时上述初始化也覆盖单格网格。

\algproblem{ALG06}{最短公共超序列}
给定字符串 $a,b$，允许空串。构造一个最短字符串 $s$，使 $a,b$ 都是 $s$ 的子序列。相同字符可以同时对应两串中的一个字符。求最短长度并恢复一个答案，给出正确性证明。
\solution 设前缀长度为 $i,j$ 时最短超序列长度为 $D_{i,j}$。边界是 $D_{i,0}=i$、$D_{0,j}=j$。若末字符相同，可让它们共用答案的最后一位：
\[
D_{i,j}=D_{i-1,j-1}+1\quad(a_i=b_j).
\]
这里字符从1编号。若末字符不同，则最后一位只能消去一个串的末字符，因此
\[
D_{i,j}=1+\min(D_{i-1,j},D_{i,j-1}).
\]
对相同末字符情形，任意超序列中两次末字符匹配都可改用较晚的那一位置，故共用最后一位不会变差；不同末字符的两个选择则穷尽了最优解。由前缀长度归纳，递推正确。

从 $(|a|,|b|)$ 回溯：相同则记录该字符并对角移动；不同则沿最小前驱记录被消去串的末字符；到边界后补入剩余前缀，最后反转记录。时间、存储均为 $O((|a|+1)(|b|+1))$。

若 $L$ 为最长公共子序列长度，则最短长度也等于 $|a|+|b|-L$：对齐两串时共用的位置形成公共子序列，最多共用 $L$ 位；反过来沿一条最长公共子序列，把两串间隙依次插入即可达到这个下界。

\algproblem{ALG07}{交织来源的计数与随机恢复}
给定字符串 $a,b,t$，其中 $|t|=|a|+|b|$。保持 $a,b$ 各自的顺序，把每个 $t$ 中的位置标为来自哪一个串。不同来源标记视为不同方案，即使字符相同。求合法方案数；若方案数非零，设计使所有合法标记等概率的随机恢复方法。
\solution 令 $C_{i,j}$ 为已经使用 $a$ 的前 $i$ 个字符、$b$ 的前 $j$ 个字符后，完成剩余 $t$ 的方案数。边界 $C_{m,n}=1$，其中 $m=|a|,n=|b|$。从后往前计算
\[
C_{i,j}=\ind_{\{i<m,\ a_{i+1}=t_{i+j+1}\}}C_{i+1,j}
       +\ind_{\{j<n,\ b_{j+1}=t_{i+j+1}\}}C_{i,j+1}.
\]
字符从1编号，越界分支贡献零。两项对应下一位置的不同来源，因此即使字符相同也应相加。总数是 $C_{0,0}$；若长度不符则直接为零。

恢复时处于 $C_{i,j}>0$ 的状态，按上述两个非负分支计数占总数的比例选择来源，再进入对应后继状态。任意完整标记的概率沿路径相乘，分子分母逐项相消，最终等于 $1/C_{0,0}$，故采样均匀。空串的唯一标记也被边界覆盖。

表含 $O((m+1)(n+1))$ 个状态，恢复恰走 $m+n$ 步。计数用精确整数；这里的状态运算次数界不把大整数加法假设为固定比特成本。只用模数计数将失去用于精确采样的权重。

\algproblem{ALG08}{数量受限的加权区间选择}
有 $n$ 个任务，任务 $i$ 占用半开区间 $[s_i,e_i)$，其中 $s_i<e_i$，收益 $w_i$ 可以为负。最多选择整数 $K\ge0$ 个互不重叠任务，允许一个任务结束时下一个立即开始。求最大总收益，要求给出 $O(n\log n+nK)$ 的算法。
\solution 按结束时间排序，记 $p(i)$ 为满足 $e_j\le s_i$ 的最大前驱下标，不存在则为0。由于 $s_i<e_i$，前驱一定排在 $i$ 之前；对有序结束时间二分即可求全部 $p(i)$，耗时 $O(n\log n)$。

令 $F(i,k)$ 是前 $i$ 个任务中至多选 $k$ 个的最好收益。边界 $F(0,k)=F(i,0)=0$，其余满足
\[
F(i,k)=\max\{F(i-1,k),\ w_i+F(p(i),k-1)\}.
\]
不选 $i$ 得到第一项；选 $i$ 时，其余所选任务必须在它开始前结束，故得到第二项。反过来第二项中的前驱最优集合与 $i$ 合并必定可行。这证明递推既是上界又可达到。

将 $K$ 截为 $\min(K,n)$，答案为 $F(n,K)$。负收益任务可被跳过，所以空集合的收益0始终是下界。完整表占 $O(n(K+1))$ 空间并支持回溯；按 $k$ 逐层计算，只求收益时可用两条长度 $n+1$ 的数组。

\algproblem{ALG09}{任务自带冷却期的日程选择}
第 $i$ 天有一项只能当天执行的任务，收益为实数 $r_i$。执行后必须休息接下来的 $L_i\ge0$ 个完整日，其中 $L_i$ 为整数；休息期间不能执行其他任务。任务可以跳过，$i=0,\ldots,T-1$。求最大总收益并恢复执行日期，要求 $O(T)$ 时间。
\solution 记 $V_i$ 为从第 $i$ 天起的最好收益，$V_T=0$。执行当天任务后，下一可执行日为
\[
q_i=\min(T,i+L_i+1).
\]
反向递推
\[
V_i=\max\{V_{i+1},r_i+V_{q_i}\},\qquad i=T-1,\ldots,0.
\]
第一项跳过当天，第二项执行并直接转入休息结束后的子问题。两种选择覆盖所有合法日程，而且拼接相应子问题的最优日程必合法，故由反向归纳得到最优性。

记录每个状态选取的分支。从0出发，跳过时前进一日，执行时输出 $i$ 并跳到 $q_i$；下标严格递增，所以恢复至多访问 $T$ 个状态。时间和空间均为 $O(T)$。若收益非正，可通过跳过得到不小于它的收益；$T=0$ 时答案0。由于 $q_i$ 的跨度各不相同，不能只保留固定数目的最近状态。

\algproblem{ALG10}{最少插入括号的修复}
给定只含左右圆括号的字符串，不能删除或改变原有字符，只能插入括号。求使整个字符串成为合法括号序列的最少插入数，并输出一种达到最少次数的修复。要求线性时间。
\solution 从左到右扫描，维护尚未配对的左括号数 $b$。遇到左括号，原样输出并令 $b$ 加1。遇到右括号：若 $b>0$，原样输出并令 $b$ 减1；若 $b=0$，先插入一个左括号，再输出该右括号，插入数加1。扫描结束后再补入 $b$ 个右括号。

输出的每个前缀左括号数不少于右括号数，末尾总数又相等，所以确为合法序列；原串的所有字符保持了顺序。

最优性可由前缀亏空证明。将空前缀也计入，记原串前缀余额的最小值为 $m\le0$，总余额为 $s$。任何修复至少插入 $L=-m$ 个左括号，否则最亏空前缀仍不合法。若插入左右括号数分别为 $\ell,r$，末尾平衡要求 $r=s+\ell$，故总插入数 $s+2\ell\ge s+2L$。算法恰插入 $L$ 个左括号，扫描后的余额为 $s+L$，再插入这么多个右括号，达到下界。时间 $O(T)$，除输出外只需常数空间。

\algproblem{ALG11}{二维地形的蓄水量}
矩形网格每格地面高度为非负实数 $h_v$，每格面积为1，相邻仅指上下左右。矩形外部可以自由排水，所有边界格都与外部相连。持续降雨至稳定后，求能够保留的总水量。给出算法并解释为什么只比较每行、每列最高点是不够的。
\solution 对每格 $v$ 定义最低逃逸水位
\[
\ell_v=\min_{\gamma:\,v\text{ 到边界}}\ \max_{u\in\gamma}h_u.
\]
路径包括起终点，故 $\ell_v\ge h_v$；某条路径的最高点决定沿它流出的门槛，而在所有逃逸路径中选最低门槛，得到稳定水位。总水量为 $\sum_v(\ell_v-h_v)$。

用最小堆计算这个“最小最大值”路径问题。所有边界格以距离 $h_v$ 入堆，内部距离初值为 $+\infty$。每次取距离最小的尚未确定格 $v$，对每个相邻格 $u$ 尝试松弛
\[
d_u\leftarrow\min\{d_u,\max(d_v,h_u)\}.
\]
忽略已过时的堆记录。由于沿路径扩展只会使最大高度不减，取出时未确定格不可能通过距离更大的路径得到更小值；这与最短路的贪心证明相同。更具体地，若有更低的逃逸路径，考察该路径上第一个未确定格，它应已由相邻的确定格获得不大的候选值，与当前最小堆顶矛盾。

网格边数为格数的常数倍，时间 $O(RC\log(RC+1))$，空间 $O(RC)$。只有一行或一列时每格都在边界，水量为零。实际逃逸路径可以转弯，因此只看直线方向的挡墙会高估蓄水。

\algproblem{ALG12}{全一正方形的最大面积与总数}
给定二元矩阵，同时求全为1的轴对齐正方形的最大面积，以及这种正方形的总个数。不同位置或边长均算不同正方形。要求 $O(RC)$ 时间、$O(C)$ 额外空间。
\solution 令 $d_{i,j}$ 为以 $(i,j)$ 作右下角的最大全一正方形边长。当前格为0时 $d_{i,j}=0$；为1时
\[
d_{i,j}=1+\min(d_{i-1,j},d_{i,j-1},d_{i-1,j-1}),
\]
越界项为0。若存在边长 $k$ 的正方形，则删掉末行、末列后，三个相邻位置都至少支持边长 $k-1$，给出必要性。反过来，三个边长 $k-1$ 的正方形与当前格的并集覆盖整个边长 $k$ 的方块，给出充分性。

最大面积是 $(\max d_{i,j})^2$。固定右下角时，可行边长恰为 $1,\ldots,d_{i,j}$，所以总数为
\[
\sum_{i,j}d_{i,j}.
\]
每个正方形有唯一右下角，因此没有重复计数。按行更新长度 $C$ 的数组时，额外保留被覆盖前的左上值；旧本列是上方，新前一列是左方。全零或空矩阵两项答案均为0。

\problem{CA01}{数据流中的加权抽样}{}
数据项按顺序到达，第 $i$ 项带正权重 $w_i$。事先不知道总项数，只允许保存一个候选项和常数个数值。设计算法，使读完任意 $t\ge1$ 项后，候选为第 $i$ 项的概率恰为 $w_i/\sum_{j=1}^t w_j$。算法每步可使用一个独立的连续均匀随机数。证明结论，并说明为什么“新项比旧候选重就替换”不正确。
\solution 维护累计权重 $W$。第一项到达时令候选为它、$W=w_1$。第 $t$ 项到达时先令 $W\leftarrow W+w_t$，再以概率 $w_t/W$ 替换候选，否则保留。一次更新只需常数次运算，保存候选标识及累计权重即可。

用归纳证明不变量。设前 $t-1$ 项累计权重为 $W_{t-1}$。旧项 $i<t$ 成为新候选，需要它原来被选中且这次未被替换；当前随机数与过去独立，故
\[
\Prob(\text{候选为 }i)
=\frac{w_i}{W_{t-1}}\left(1-\frac{w_t}{W_t}\right)
=\frac{w_i}{W_t}.
\]
新项成为候选的概率则直接为 $w_t/W_t$。归纳起点 $t=1$ 显然成立，因此对每个确定的 $t$ 都满足要求；结束长度由外部给定且不依赖抽样结果时，结论也成立。

例如权重依次为 $1,2,3$，第二步选第二项的概率为 $2/3$，第三步替换概率为 $1/2$，最终三项概率为 $1/6,1/3,1/2$。按候选权重大小作确定性替换，在严格递增的权重流中总会保留最后一项，无法产生所需分布。

算法对固定权重序列成立；若权重随机，但与抽样所用随机数独立，可以先条件于整条权重序列应用同一证明。若后续权重由已公开的候选结果自适应决定，不能直接沿用这个条件化结论。空间界按保存的数值个数计算，精确累计权重所需的位数可能增长。
\insight 保存的是满足概率不变量的代表项。新项应与全部历史权重竞争，而不是只与当前候选的权重比较。

\problem{CA02}{两个有序数组的中位数}{}
给定两个非递减实数数组 $A,B$，长度分别为 $m,n$，且 $m+n>0$。允许随机访问，不允许先合并全部元素。求合并样本的中位数：总数为偶数时取中间两数的平均值。设计 $O(\log(\min\{m,n\}+2))$ 时间、$O(1)$ 额外空间的算法，处理空数组、重复值和极不相同的数组长度。
\solution 必要时交换两数组，使 $m\le n$。若 $m=0$，直接读取 $B$ 的中央一项或两项。以下设 $m>0$，令
\[
h=\left\lfloor\frac{m+n+1}{2}\right\rfloor.
\]
从 $A$ 的前部取 $i$ 项、从 $B$ 的前部取 $j=h-i$ 项作为左半部。合法搜索区间为
\[
\max(0,h-n)\le i\le\min(m,h).
\]
定义分割两侧的边界元素
\[
\begin{array}{ll}
L_A=A_{i-1}\ (i>0),& R_A=A_i\ (i<m),\\
L_B=B_{j-1}\ (j>0),& R_B=B_j\ (j<n).
\end{array}
\]
缺失的左边界视为 $-\infty$，缺失的右边界视为 $+\infty$；这些哨兵仅用于比较。

正确分割的充要条件是 $L_A\le R_B$ 且 $L_B\le R_A$。因为数组各自有序，这两个交叉比较就保证整个左半部不大于整个右半部。若 $L_A>R_B$，说明从 $A$ 取多了，应减小 $i$；随着 $i$ 增大，$L_A$ 不减而 $R_B$ 不增，故所有更大的 $i$ 都不可能正确。若 $L_B>R_A$，同理应增大 $i$。据此在合法整数区间二分。

正确分割必存在：将全部元素概念上排序并对相等值固定任一稳定顺序，前 $h$ 项在每个原数组内都是前缀。二分的每次排除都保留至少一个正确分割，因此必终止。找到后，若 $m+n$ 为奇数，答案为 $\max(L_A,L_B)$；若为偶数，答案为
\[
\frac{\max(L_A,L_B)+\min(R_A,R_B)}2.
\]
总数为偶数且非零时两半均非空，因此最终平均的都是有限真实元素。重复值允许多个正确分割，但它们给出的中位数相同。搜索区间至多含 $m+1$ 个整数，得到所需时间和空间界。
\insight 二分搜索的对象是左半部应从较短数组取多少项。证明方向单调后，不必实际合并任何元素。

\problem{CA03}{逆序对与最少相邻交换}{}
实数数组 $a_1,\ldots,a_n$ 可有重复值。逆序对定义为 $i<j$ 且 $a_i>a_j$。设计 $O(n\log n)$ 时间算法计数，并证明将数组变成非递减顺序所需的最少相邻交换次数恰等于逆序对数。计算 $(3,1,2,1)$ 的结果，说明相等元素的处理。
\solution 用归并排序同时计算逆序对。先分别递归处理左右两半，累计内部逆序对；再将已排序的两半合并。若当前左值不大于右值，就输出左值。否则输出右值，且左半部所有尚未输出的值都严格大于该右值，故一次增加“左边剩余元素数”个跨半区逆序对。每对跨半区元素只在其右元素先被输出时计数一次，因此不遗漏、不重复。

每层合并总工作量为 $O(n)$，层数为 $O(\log n)$，额外缓冲区为 $O(n)$。空数组或单元素数组的计数为零。相等时优先输出左值，避免把相等元素误计为逆序；计数变量须能容纳 $n(n-1)/2$。

再证明交换数的最优性。交换一对相邻元素时，它们与其他元素形成的逆序对总数不变，因为其他元素位于这两项的同一侧；只有交换的这对关系可能改变。因此一次交换使总逆序数至多下降一。排序后的逆序数为零，所以任何方案至少需要初始逆序数次交换。

若数组尚未有序，必有相邻的严格下降对，否则所有相邻关系非递减就已保证整体有序。交换这样的下降对，逆序数恰减少一。重复操作，恰在初始逆序数步后到达有序数组，达到下界。例中逆序对为
\[
(1,2),(1,3),(1,4),(3,4),
\]
总数为 $4$，故最少相邻交换次数也是 $4$。
\insight 分治使每一对元素有唯一的计数时机；逆序数又是证明交换操作最优性的势函数。


\section{练习}
\noindent\textbf{要点与方法.}\enspace
本组练习检验摘要的可恢复性、窗口的合法性和候选删除规则。方差计算先区分总体方差与样本方差的分母；区间问题先明确下标、长度和并列规则。给出一个可行区间只证明了可行性，还应排除更优区间；说明线性时间时，要累计整个扫描过程中每个候选被处理的次数。
\par\smallskip


\exercise{AG4-1}{恢复被合并的数据组}
两组互不重叠的数据合并后共有5个数，均值4，中心平方和26。已知其中 A 组为 $(1,4,7)$。不读取 B 组原始数据，恢复其个数、均值与中心平方和，并求合并样本的总体方差和样本方差。能否唯一恢复 B 组的两个数？

\exercise{AG4-2}{最长稳定区间的端点}
序列为 $(4,2,3,5,8,7,6)$。找出最大值与最小值之差不超过3的最长连续区间，采用从0开始的下标，长度并列时取最早起点。给出结果，并用一个简短的独立论证排除更长区间。

\exercise{AG4-3}{期限约束排除了哪个峰值}
价格依次为 $(120,80,100,70,90)$。分别求持续时间不超过2与不设上限时的最大回撤，并给出从0开始的区间下标。解释为何全历史最高价不能直接用于前一种计算。

\exercise{AG4-4}{达到四分之一回撤的最短时间}
价格依次为 $(100,94,80,79,60)$，阈值为 $d=1/4$。求达到阈值的最短回撤区间，给出下标。验证没有更短的可行区间，并说明终点扫描后为何能永久删除一个已经达标的起点。
